Diante disto, esta revisão de escopo tem por objetivo validar epistemologicamente os Saberes e Sistemas Agrícolas Tradicionais mediante a aplicação de algoritmos de Aprendizado de Máquina para decodificar e quantificar processos de inovação tácita. A hipótese central que orienta esta investigação postula que os algoritmos de Aprendizado de Máquina, ao atuarem como decodificadores de capacidades dinâmicas, transcendem a função de mapeamento para operar a validação epistemológica do conhecimento tácito. Ao converter saberes e sistemas das comunidades em parâmetros auditáveis, a modelagem computacional estabelece a ponte necessária entre a sabedoria ancestral e os mercados de valor agregado, assegurando que a proteção dos ativos intangíveis e a integridade ecológica desses sistemas sejam reconhecidas não como externalidades culturais, mas como vetores de competitividade e resiliência no Antropoceno.
